\documentclass[11pt]{article}
\usepackage{amsmath,amssymb,amsthm}
\usepackage[margin=1.2in]{geometry}
\newtheorem{theorem}{Theorem}
\newtheorem{lemma}[theorem]{Lemma}
\newtheorem{corollary}[theorem]{Corollary}
\newtheorem{proposition}[theorem]{Proposition}
\newtheorem{remark}[theorem]{Remark}
\title{The turn-division ledger:\\ what happens when the half-turn is replaced by a
$1/q$-turn}
\author{Samuel Lavery}
\date{Draft --- \today}
\begin{document}
\maketitle

\begin{abstract}
Every carry criterion used so far in this campaign reads the phase $\{n/p^{j}\}$ against
$\tfrac12$.  That threshold is not intrinsic: it is an artifact of $\binom{2n}{n}$ splitting
$2n$ into \emph{two} equal parts.  We replace the half-turn by a division of the circle into
$q$ equal sectors, prove the corresponding ledger identity, and re-derive every half-turn
result at general $q$.  Three consequences.  (i) The antipodal law survives verbatim at every
$q$, with threshold $\max(qk,\sqrt{qn})$.  (ii) The $q$-analogue of Erd\H{o}s \#396 is
\emph{settled outright} whenever $q\ge k+1$, by an explicit construction; so the difficulty of
\#396 is the difficulty of its extreme case $q=2$, not of the divisibility itself.  (iii)
Exactly one rigidity survives scrutiny: the registration offset has mean $\tfrac1{2p}$ at
\emph{every} division, and the half-turn is distinguished solely by its offset having zero
\emph{variance}, which happens iff $\varphi(q)=1$.  Two further candidate rigidities are
disposed of --- the ceiling maximum is monotonicity in $q$, and the $\pi/3$ cell is an
artifact of the chart.
\end{abstract}

\section{The ledger at a $1/q$-turn}

Throughout $q\ge2$ is an integer, $p$ a prime, $s_p$ the base-$p$ digit sum, $\{x\}$ the
fractional part.  Write
\[
  M_q(n)\ :=\ \frac{(qn)!}{(n!)^{q}},
\]
so $M_2(n)=\binom{2n}{n}$.  $M_q(n)$ is the multinomial coefficient counting the ordered
partitions of a $qn$-set into $q$ blocks of size $n$, hence a positive integer.

\begin{theorem}[turn-division ledger]\label{thm:ledger}
For every prime $p$ and every $n\ge0$,
\[
  v_p\bigl(M_q(n)\bigr)\ =\ \sum_{j\ge1}\Bigl\lfloor\,q\,\bigl\{\,n/p^{j}\,\bigr\}\,\Bigr\rfloor .
\]
\end{theorem}

\begin{proof}
By Legendre, $v_p(M_q(n))=\sum_{j\ge1}\bigl(\lfloor qn/p^{j}\rfloor-q\lfloor
n/p^{j}\rfloor\bigr)$.  Writing $x=n/p^{j}=\lfloor x\rfloor+\{x\}$ gives
$\lfloor qx\rfloor=q\lfloor x\rfloor+\lfloor q\{x\}\rfloor$, and the claim follows termwise.
\end{proof}

\emph{Verified}: $5040$ triples $(q,p,n)$ with $2\le q\le8$ --- no mismatch.

At $q=2$ each summand is the half-turn indicator $\mathbf 1[\{n/p^{j}\}\ge\frac12]$, taking
only the values $0,1$.  For $q\ge3$ the summand is the \emph{index of the sector}: the circle
is cut into $q$ equal arcs and the ledger records which arc the phase occupies, a value in
$\{0,1,\dots,q-1\}$.  This is the only change, and every statement below is obtained by
tracking it.

\begin{corollary}[the $q$-carrier]\label{cor:carrier}
$p\nmid M_q(n)$ if and only if every base-$p$ digit of $n$ is $<p/q$.  Write
$G_p^{(q)}$ for this set; $G_p^{(2)}$ is the carrier of \textup{\#377}.
\end{corollary}

\begin{proof}
The $j$-th summand vanishes iff $\{n/p^{j}\}<1/q$, i.e.\ iff $n\bmod p^{j}<p^{j}/q$.  Taking
$j=1$ gives the condition on the last digit; applying it to $\lfloor n/p^{j-1}\rfloor$ for
each $j$ gives it on every digit, and conversely.
\end{proof}

\emph{Verified}: $28\,000$ triples $(q,p,n)$, $2\le q\le8$, $p<30$, $n<400$ --- no mismatch.

The number of admissible digits is $\lceil p/q\rceil$, so the codimension (box-counting cost)
of confining a rail to $G_p^{(q)}$ is
\[
  \beta_p^{(q)}\ :=\ \frac{\log\bigl(p/\lceil p/q\rceil\bigr)}{\log p} .
\]
A rail is \emph{degenerate} when $\lceil p/q\rceil=1$, i.e.\ when $p\le q$: then
$G_p^{(q)}=\{0\}$ and $\beta_p^{(q)}=1$.  Degenerate rails must be excluded from any budget
argument; including them returns the spurious optimum $\{2\}$ at every $q\ge3$.

\section{The antipodal law at general $q$}

\begin{theorem}[general antipodal law]\label{thm:anti}
Let $p>\max\bigl(qk,\ \sqrt{qn}\bigr)$ be prime and put $\rho=n\bmod p$.  Then
$v_p(M_q(n))=\lfloor q\rho/p\rfloor$, and
\begin{enumerate}
\item[(a)] if $p\mid n-i$ for some $0\le i\le k$ then $v_p(M_q(n))=0$, so $p$ cannot divide
      $M_q(n)$: \emph{extinction};
\item[(b)] if $p\mid n+i$ for some $1\le i\le k$ then $v_p(n+i)=1\le v_p(M_q(n))$:
      \emph{automatic}.
\end{enumerate}
\end{theorem}

\begin{proof}
For $j\ge2$ we have $p^{j}\ge p^{2}>qn\ge n$, so $\{n/p^{j}\}=n/p^{j}$ and
$q\{n/p^{j}\}=qn/p^{j}<1$; that level contributes $0$.  Only $j=1$ survives in
Theorem~\ref{thm:ledger}, and there $\{n/p\}=\rho/p$, giving
$v_p(M_q(n))=\lfloor q\rho/p\rfloor$.
(a) $p\mid n-i$ with $0\le i\le k$ forces $\rho=i$, and $i\le k<p/q$ because $p>qk$; so
$q\rho/p<1$ and the entry is $0$.  (b) $p\mid n+i$ with $1\le i\le k$ forces $\rho=p-i$, and
$qi\le qk<p$, so $q\rho/p=q-qi/p>q-1\ge1$ and the entry is at least $1$; finally
$p^{2}>qn\ge n+i$ gives $v_p(n+i)=1$.
\end{proof}

\emph{Verified}: for $q=2,3,4,5,6$ and $k\le6$, $n<900$: $66\,896$ negative-shift incidences,
all with $v_p(M_q)=0$, and $50\,165$ positive-shift incidences, all with
$v_p(n+i)=1\le v_p(M_q)$ --- no failure in either direction at any $q$.

Two sharpenings, both of which are invisible at $q=2$.

\begin{proposition}[the positive shift saturates]\label{prop:sat}
In case \textup{(b)} the value is exact: $v_p(M_q(n))=q-1$, the maximal sector.  Hence
$\prod_{1\le i\le k}(n+i)^{\,q-1}$ divides $M_q(n)$ as far as the primes
$p>\max(qk,\sqrt{q(n+k)})$ are concerned.
\end{proposition}

\begin{proof}
$\rho=p-i$ with $1\le i\le k$ and $qk<p$ give $q-1<q-qi/p=q\rho/p<q$, so
$\lfloor q\rho/p\rfloor=q-1$.
\end{proof}

\emph{Verified}: $44\,688$ incidences over $2\le q\le7$, $k\le6$, $n<700$ --- no failure.  So
coarsening the division strengthens the positive shift (to the $(q-1)$-th power) while leaving
the negative shift exactly as dead; at $q=2$ the exponent is $1$ and the asymmetry cannot be
seen.

\begin{theorem}[$a$-fold extinction]\label{thm:afold}
Let $a\ge1$, let $p>qk$ be prime with $p^{a+1}>qn$, and suppose $p^{a}\mid n-i$ for some
$0\le i\le k$.  Then $v_p(M_q(n))=0$.
\end{theorem}

\begin{proof}
$n\equiv i\pmod{p^{a}}$ with $i\le k<p$ gives $n\bmod p^{j}=i$ for every $j\le a$, so
$\lfloor q\{n/p^{j}\}\rfloor=\lfloor qi/p^{j}\rfloor\le\lfloor qk/p\rfloor=0$ at each such
level; and for $j>a$ we have $p^{j}\ge p^{a+1}>qn$, so those levels vanish too.
\end{proof}

\emph{Verified}: $79\,479$ incidences over $2\le q\le7$, $k\le6$, $n<900$ --- no failure.

Theorem~\ref{thm:afold} is a \emph{tower} of conditions, not one.  Taking $a=1$ recovers
Theorem~\ref{thm:anti}(a) and forces the window to be $(qn)^{1/2}$-smooth; $a=2$ forces its
square-full part to be $(qn)^{1/3}$-smooth; $a=3$ its cube-full part to be $(qn)^{1/4}$-smooth;
and so on.  At $q=2$ with multiplicity one this is invisible, which is why the \#396 note
recorded only the first floor.

The tower is what kills the obvious constructions, and it does so uniformly in $q$.  The two
parametric families that \emph{do} force a smooth window --- $n=x^{2}$ at $k=1$, with window
$\{x^{2},(x-1)(x+1)\}$, and the Pell family $y^{2}-2x^{2}=-1$, $n=2x^{2}=y^{2}+1$ at $k=2$ ---
are built out of squares, and squares are exactly what produces $v_p\ge2$ at $p\approx\sqrt n$,
where $p^{3}>qn$.  Concretely $n=1682=2\cdot29^{2}$ has $v_{29}(n)=2$ but
$v_{29}(M_q(n))=0$ for every $q$ from $2$ to $14$; all eleven Pell terms up to
$n\approx10^{17}$ die this way at every $q$ tested.  The $n=x^{2}$ family survives only when
$x$ is itself $(qx^{2})^{1/3}$-smooth.  And the Chinese-remainder construction has no reach
past $k=1$ at any $q$: forcing the window smooth needs $k+1$ pairwise-coprime moduli of size
about $\sqrt n$, whose product $n^{(k+1)/2}$ must not exceed $n$.

So the phase picture is $q$-independent in form: the divisor condition for a \emph{negative}
shift pins $\rho$ to the bottom arc $[0,k]$, which is disjoint from the carry arc
$[p/q,p)$ as soon as $k<p/q$; a \emph{positive} shift pins $\rho$ to the top arc
$[p-k,p)$, which lies inside the carry arc.  Only the threshold moves, from
$\max(2k,\sqrt{2n})$ to $\max(qk,\sqrt{qn})$.

\section{The non-divisor band at general $q$}

Two earlier notes isolate, at $q=2$, the primes that fail to divide $\binom{2n}{n}$: the band
$(2n/3,n]$ and, more generally, the interval criterion $n\in G_p\iff p(2m+1)\ge2n+1$ and
$p\ge2m+1$, with $m=\lfloor n/p\rfloor$.  Both are the $q=2$ case of one statement.

\begin{theorem}[non-divisor band]\label{thm:band}
Let $p>\sqrt{qn}$ be prime and put $m=\lfloor n/p\rfloor\ge1$.  Then
\[
  p\nmid M_q(n)
  \iff
  p\,(qm+1)>qn
  \iff
  p\ \in\ I^{(q)}_m(n)\ :=\ \Bigl(\ \frac{qn}{qm+1},\ \frac nm\ \Bigr] .
\]
The band has width $\dfrac{n}{m\,(qm+1)}$, and the top band $m=1$ is
$\bigl(\tfrac{q}{q+1}n,\ n\bigr]$.
\end{theorem}

\begin{proof}
As in Theorem~\ref{thm:anti}, $p^{2}>qn$ leaves only the level $j=1$, so
$v_p(M_q(n))=\lfloor q\{n/p\}\rfloor$.  This vanishes iff $\{n/p\}<1/q$, i.e.\ iff
$n/p-m<1/q$, i.e.\ iff $p>n/(m+\tfrac1q)=qn/(qm+1)$.  The upper endpoint is
$m=\lfloor n/p\rfloor$, i.e.\ $p\le n/m$.  Subtracting the endpoints gives the width.
\end{proof}

\emph{Verified}: $1\,020\,921$ triples $(q,n,p)$ with $2\le q\le7$, $n<1500$,
$\sqrt{qn}<p\le n$ --- no mismatch.

At $q=2$ the condition $p(2m+1)>2n$ is, for integers, exactly $p(2m+1)\ge2n+1$, and the
side condition $p>\sqrt{2n}$ is the interval criterion's $p\ge2m+1$; the top band is
$(2n/3,n]$.  At $q=3$ the top band is $(3n/4,n]$, at $q=4$ it is $(4n/5,n]$: the coarser the
division, the \emph{narrower} the non-divisor bands --- consistent with
Proposition~\ref{prop:ceiling}, since a coarser division makes carries easier and
non-divisors rarer.

\begin{corollary}[criterion at general $q$]\label{cor:crit}
For $n\ge k+1$,
\[
  \prod_{0\le i\le k}(n-i)\ \Bigm|\ M_q(n)
  \iff
  (q+1)\,s_p(n)\ \ge\ (k+1)+s_p(n-k-1)+s_p(qn)\ \ \text{for every prime }p .
\]
\end{corollary}

\begin{proof}
By Legendre the left side of the valuation inequality is
$\bigl((k+1)-s_p(n)+s_p(n-k-1)\bigr)/(p-1)$ and the right side is
$\bigl(q\,s_p(n)-s_p(qn)\bigr)/(p-1)$; clear denominators and collect.
\end{proof}

At $q=2$ this is the criterion $3s_p(n)\ge(k+1)+s_p(n-k-1)+s_p(2n)$ of the \#396 note.
\emph{Verified}: reproduces the least solutions $1,2,2480,8178,45153$ for $k=0,\dots,4$ at
$q=2$.

\section{The $q$-analogue of \#396 is settled for $q\ge k+1$}

\begin{theorem}\label{thm:qge}
Let $k\ge0$ and $q\ge k+1$.  Then $n=k+1$ satisfies
$\displaystyle\prod_{0\le i\le k}(n-i)\ \Bigm|\ M_q(n)$.
\end{theorem}

\begin{proof}
Put $m=k+1$, so $n=m$ and $\prod_{0\le i\le k}(n-i)=m(m-1)\cdots1=m!$.  The number of ways to
partition a set of $qm$ elements into $q$ \emph{unordered} blocks of size $m$ is
$\frac{(qm)!}{(m!)^{q}\,q!}$, an integer; hence $q!$ divides $M_q(m)=\frac{(qm)!}{(m!)^{q}}$.
Since $m\le q$ we have $m!\mid q!$, so $m!$ divides $M_q(m)$.
\end{proof}

\emph{Verified}: $104$ pairs $(q,k)$ with $2\le q\le14$, $0\le k<q$, by exact integer
arithmetic --- no failure.

The hypothesis $q\ge k+1$ is not the exact condition, and the exact condition explains the
apparent irregularity at the boundary.

\begin{proposition}[when the construction $n=k+1$ works]\label{prop:exact}
Put $m=k+1$.  Then $\prod_{0\le i\le k}(m-i)\mid M_q(m)$ if and only if
\[
  (q+1)\,s_\ell(m)\ \ge\ m+s_\ell(qm)\qquad\text{for every prime }\ell .
\]
If moreover $m$ is a power of the prime $\ell$, the condition at that $\ell$ reads
$q-s_\ell(q)\ge k$; in particular it fails whenever $q\le k$.
\end{proposition}

\begin{proof}
Substitute $n=m=k+1$ into Corollary~\ref{cor:crit}, where $n-k-1=0$ and $s_\ell(0)=0$.  If
$m=\ell^{j}$ then $s_\ell(m)=1$ and $s_\ell(qm)=s_\ell(q\ell^{j})=s_\ell(q)$, so the
inequality becomes $q+1\ge m+s_\ell(q)$, i.e.\ $q-s_\ell(q)\ge m-1=k$.  Since $s_\ell(q)\ge1$
for $q\ge1$ this forces $q\ge k+1$.
\end{proof}

\begin{corollary}[the boundary $k=q$]\label{cor:boundary}
At $k=q$ the construction $n=k+1$ fails whenever $q+1$ is a prime power.
\end{corollary}

\begin{proof}
Take $m=q+1=\ell^{j}$ in Proposition~\ref{prop:exact}: the condition at $\ell$ is
$q-s_\ell(q)\ge k=q$, i.e.\ $s_\ell(q)\le0$, which is false.
\end{proof}

Corollary~\ref{cor:boundary} has a phrasing free of $q$ and $k$ altogether.  Putting
$n=q+1$, the criterion at $k=q$ reads $n(s_\ell(n)-1)\ge s_\ell(n(n-1))$ for every prime
$\ell$, and the divisibility being tested is
\[
  (n!)^{\,n}\ \Bigm|\ \bigl(n(n-1)\bigr)! \qquad\text{iff}\qquad n\ \text{is not a prime power.}
\]
\emph{Verified} by exact factorial arithmetic for $2\le n\le25$ --- no mismatch.  This is a
clean standalone statement and I have \emph{not} run a literature check on it; it should be
looked up before anyone calls it new.  (The neighbouring fact $(n!)^{\,n+1}\mid(n^{2})!$ is
classical, being the count of partitions of an $n^{2}$-set into $n$ unordered $n$-blocks.)

\emph{Measured}: the converse also holds throughout $2\le q\le259$ --- at $k=q$ the
construction succeeds for exactly those $q$ with $q+1$ \emph{not} a prime power, i.e.\
$q+1\in\{6,10,12,14,15,18,20,21,22,24,26,\dots\}$, and fails for exactly those with $q+1$ a
prime power, $q+1\in\{3,4,5,7,8,9,11,13,16,17,19,23,25,27,\dots\}$; $258$ values of $q$, no
mismatch, and an independent check to $q=39$ agrees.  The converse is not proved here: the
forward direction is Corollary~\ref{cor:boundary}, but showing that $s_\ell(m)\ge2$ for all
$\ell$ suffices needs more than the crude digit bound, since $(\ell-1)(2\log_\ell m+1)$ can
exceed $m$ --- and the standard submultiplicative bound $s_\ell(ab)\le s_\ell(a)s_\ell(b)$ is
exactly tight at $\ell=m/2$, so it does not close the gap either.  The measured margin
$\min_\ell\bigl[m(s_\ell(m)-1)-s_\ell(m(m-1))\bigr]$ over all non-prime-power $m<400$ is
always $\ge0$, with minimum $+2$ attained at $m=6$.

The law is specific to the boundary $k=q$: at $k=q+1$ it fails, with $26$ mismatches in $118$
values, the first at $q=4$.  It is not a general prime-power principle.

\begin{corollary}
For every $k$ there is a division $q$ and an $n$ with $\prod_{0\le i\le k}(n-i)\mid M_q(n)$;
one may take $q=k+1$ and $n=k+1$.
\end{corollary}

Table of the least $n$ with $\prod_{0\le i\le k}(n-i)\mid M_q(n)$, computed from
Corollary~\ref{cor:crit} (entries $>4\times10^{5}$ were not found in the search range):
\[
\begin{array}{r|rrrrrrr}
 k & q{=}2 & q{=}3 & q{=}4 & q{=}5 & q{=}6 & q{=}8 & q{=}12\\\hline
 0 & 1 & 1 & 1 & 1 & 1 & 1 & 1\\
 1 & 2 & 2 & 2 & 2 & 2 & 2 & 2\\
 2 & 2480 & 3 & 3 & 3 & 3 & 3 & 3\\
 3 & 8178 & 496 & 4 & 4 & 4 & 4 & 4\\
 4 & 45153 & 13068 & 496 & 5 & 5 & 5 & 5\\
 5 & \cdot & 44255 & 2627 & 6 & 6 & 6 & 6\\
 6 & \cdot & 48512 & 2628 & 2628 & 10 & 7 & 7\\
 7 & \cdot & 48513 & 48513 & 48513 & 10 & 8 & 8\\
 8 & \cdot & \cdot & \cdot & \cdot & 10 & 10 & 9\\
 9 & \cdot & \cdot & \cdot & \cdot & 10 & 10 & 10\\
10 & \cdot & \cdot & \cdot & \cdot & \cdot & 15 & 11
\end{array}
\]
The diagonal band of entries $n=k+1$ is exactly Theorem~\ref{thm:qge}, and it terminates at
$k=q-1$ in every column.  Read across a row: the difficulty of \#396 is monotone in the
\emph{coarseness} of the division, and the half-turn is the worst case because $M_2$ is the
smallest of the $M_q$ and therefore has the least carry budget to spend.  Monotonicity in $q$
is provable, not just observed: each sector index $\lfloor q\{n/p^{j}\}\rfloor$ is
nondecreasing in $q$, so the solution set is upward-closed in $q$ and $n_{\min}(q,k)$ is
nonincreasing in $q$.

\emph{The end of the diagonal is a cliff, not a wall.}  It would be wrong to read the table as
saying solutions exist only for $k<q$: they exist well beyond, they are merely large.  An
independent search extends the rows to
$n_{\min}(2,\cdot)=\dots,3648841,\,7979090,\,101130029,\,339949252$ for $k=5,\dots,8$;
$n_{\min}(3,11)=321903043$; $n_{\min}(6,14)=216715180$.  What happens at the end of the
diagonal is a jump: at $q=6$ the value sits at $10$ for $k=6,7,8$ and then leaps to $376251$
at $k=9$, a factor of $3.8\times10^{4}$ in one step, exactly where $n\le qk$ fails.  Fitted
growth of $n_{\min}$ over the last six $k$ in each row: $9.05^{k}$ at $q=2$, $5.50^{k}$ at
$q=3$, $4.79^{k}$ at $q=5$, $3.55^{k}$ at $q=6$.  So coarsening the division changes the base
of the exponential; it does not remove it.  Theorem~\ref{thm:qge} bounds \emph{the reach of
one construction}, not the existence of solutions.

\section{In what sense the half-turn is special: three candidates, one survivor}

\begin{theorem}[registration offset: DC and AC]\label{thm:phi}
Write $\lceil p/q\rceil/p=\frac1q+\frac{r}{qp}$ with $r=(-p)\bmod q$.  As $p$ ranges over the
primes exceeding $q$, the value of $r$ runs over $R_q:=\{(-a)\bmod q:\gcd(a,q)=1\}$, a set of
size $\varphi(q)$, equidistributed by Dirichlet.  Then:
\begin{enumerate}
\item[\rm(DC)] the mean of $r$ over $R_q$ is $q/2$ for \emph{every} $q\ge2$, so the mean
      registration offset is $\dfrac1{2p}$ at every division;
\item[\rm(AC)] the variance of $r$ over $R_q$ vanishes if and only if $\varphi(q)=1$, that is
      $q\in\{1,2\}$.
\end{enumerate}
Hence the half-turn is the unique division whose registration offset has \emph{zero AC
content}.
\end{theorem}

\begin{proof}
$q\lceil p/q\rceil=p+r$ gives the displayed identity, and $\gcd(p,q)=1$ for $p>q$ prime, so by
Dirichlet $r$ takes exactly the values $R_q$.  For (DC): $a\mapsto q-a$ maps
$(\mathbb Z/q)^{\times}$ to itself and has no fixed point when $q\ge3$ (a fixed point needs
$2a=q$, whence $\gcd(a,q)=a=q/2>1$), so $R_q$ is a disjoint union of pairs $\{r,q-r\}$ and its
mean is $q/2$; for $q=2$, $R_2=\{1\}$ and the mean is $1=q/2$.  For (AC): a multiset with mean
$q/2$ has zero variance iff it is the singleton $\{q/2\}$, iff $|R_q|=\varphi(q)=1$.
\end{proof}

\emph{Verified}: exact $(\mathbb E[r],\operatorname{Var}[r])$ over $R_q$ for $2\le q\le12$ ---
$q=2:(1,0)$, $q=3:(\tfrac32,\tfrac14)$, $q=4:(2,1)$, $q=5:(\tfrac52,\tfrac54)$, $q=6:(3,4)$,
$q=8:(4,5)$, $q=12:(6,13)$ --- mean equal to $q/2$ in every case and variance zero only at
$q=2$.

This correction matters, and it reverses an earlier reading of mine.  It is \emph{not} true
that the offset $\tfrac1{2p}$ is the half-turn's signature: by (DC) that offset is the common
mode of every division.  What the half-turn alone has is a \emph{fluctuation-free} offset, so
that the offsets can be summed rail by rail with no error term; for $q\ge3$ the fluctuation is
carried by the odd Dirichlet characters mod $q$ and no single ledger entry serves all rails.

There is an intrinsic, chart-free companion to (AC).  At $q=2$ the nonzero admissible digits
are $\{1,\dots,\tfrac{p-1}2\}$, which is exactly a fundamental domain for $x\mapsto-x$ on
$(\mathbb Z/p)^{\times}$: each pair $\{x,-x\}$ meets it once.  So the admissible set is not
merely an interval, it is a transversal of $(\mathbb Z/p)^{\times}/\{\pm1\}$, for every odd
$p$.  Both (AC) and this statement are avatars of the same involution $x\mapsto-x$:
$\varphi(q)=1$ says $-1\equiv1\pmod q$, and $\{\pm1\}$ is the subgroup whose cosets an
interval can cut out.  For $q\ge3$ no such group-theoretic description of the admissible
interval exists except sporadically.

\begin{theorem}[the limiting cost angle]\label{thm:niven}
Define $\theta_p^{(q)}\in[0,\pi/2)$ by $\cos\theta_p^{(q)}=\lceil p/q\rceil/p$, so that
$\beta_p^{(q)}=\log\sec\theta_p^{(q)}/\log p$, and let
$\theta_\infty^{(q)}=\lim_{p\to\infty}\theta_p^{(q)}=\arccos(1/q)$.  Then
$\theta_\infty^{(q)}$ is a rational multiple of $\pi$ if and only if $q\in\{1,2\}$, and for
$q=2$ it equals $\pi/3$.
\end{theorem}

\begin{proof}
$\cos\theta_\infty^{(q)}=1/q$ is rational, so by Niven's theorem --- the only rational values
of $\cos\theta$ at $\theta$ a rational multiple of $\pi$ are $0,\pm\frac12,\pm1$ --- we need
$1/q\in\{0,\tfrac12,1\}$, i.e.\ $q\in\{1,2\}$.  At $q=2$, $\arccos\frac12=\pi/3$.
\end{proof}

The theorem is true, and its content is nil.  The deduction uses only
``$1/q=\tfrac12\iff q=2$'': any rational-valued family meets the five-element Niven set exactly
where it equals one of five numbers, and Niven merely re-expresses that as a statement about
angles.  The apparent payoff --- that the half-turn lands on the $\mu_6$ cell --- is an
artifact of the chart, for four reasons.

\emph{(i) Other equally natural charts give other cells, or none.}  Taking
$\sin\theta=1/q$ also selects $q=2$ but at $\theta=\pi/6$, the $\mu_{12}$ cell, not $\mu_6$.
Taking $\tan\theta=1/q$ selects nothing, since Niven's tangent set is $\{0,\pm1\}$ and forces
$q=1$.  The density $1/q$ read bare supports no lattice statement at all.  And the
codimension $\beta_p^{(q)}$ is \emph{irrational for every nondegenerate pair} $(p,q)$: if
$\log(p/a)/\log p=u/v$ then $a^{v}=p^{v-u}$, so $a$ is a power of $p$, and $1\le a\le p$
forces $a=1$ ($\beta=1$) or $a=p$ ($\beta=0$).  So $\beta$ never lands on any lattice, at any
$q$.  Perturbing the family decouples the two conditions further: $\cos\theta=1/(q-1)$ selects
$\{2,3\}$ with $q=3$ on $\mu_6$, and $\cos\theta=(q-1)/(q+1)$ selects $\{3\}$ alone.

\emph{(ii) The object supplies its own angle, and it distinguishes nothing.}  The admissible
digits are precisely the $p$-th roots of unity lying in the arc $[0,2\pi/q)$.  That arc has
width $2\pi/q$, a rational multiple of $\pi$ for \emph{every} $q$; in the construction's own
angular coordinate each division sits on $\mu_q$ and no division is special.

\emph{(iii) The exact cell is occupied only by the trivial rail.}  Since
$\delta=\lceil p/q\rceil/p\in(0,1]$, the values $0,-\tfrac12,-1$ are unreachable; $\delta=1$
iff $q=1$; and $\delta=\tfrac12$ needs $2\lceil p/q\rceil=p$, hence $p$ even, hence $p=2$.  So
the complete list of exact Niven landings is: $q=1$ with every $p$, and $p=2$ with
\emph{every} $q\ge2$, where $\delta=\tfrac12$ and $\theta=\pi/3$ exactly.  The sole exact
$\mu_6$ landing in the whole family is therefore the degenerate rail $p=2$, and it occurs at
every division.

\emph{(iv) The $6$ has no preimage in the object.}  At $q=2$ the construction produces $2$
everywhere --- arc width $\pi$, stabiliser $\{\pm1\}$, offset modulus $2$.  The $6$ enters
only as the order of the root of unity attached to $\arccos\tfrac12$, i.e.\ from the chart.
A cell whose order cannot be exhibited anywhere in the object is not structure.

What survives is the bare selection: $\arccos(1/q)/\pi\notin\mathbb Q$ for every $q\ge3$.  The
cell does not.

\begin{proposition}[the ceiling collapses]\label{prop:ceiling}
Let $\Lambda(q)=\max\bigl\{\sum_{p\in S}1/p\ :\ S\ \text{a set of nondegenerate rails with}\
\sum_{p\in S}\beta_p^{(q)}<1\bigr\}$.  Over rails $p<300$ and $|S|\le5$ the computed values
are
\[
\begin{array}{r|cccccc}
q & 2 & 3 & 4 & 5 & 6 & 7\\\hline
S & \{3,5,7\} & \{5,11\} & \{5,23\} & \{7,89\} & \{7,139\} & \{11,827\}\\
\Lambda & 0.676190 & 0.290909 & 0.243478 & 0.154093 & 0.150051 & 0.092118\\[2pt]\hline
q & 8 & 9 & 10 & 11 & 12 & \\\hline
S & \{11,1321\} & \{11,1993\} & \{11,2851\} & \{13,7109\} & \{13,9817\} & \\
\Lambda & 0.091666 & 0.091411 & 0.091260 & 0.077064 & 0.077025 &
\end{array}
\]
Every entry is confirmed by an independent branch-and-bound search in ratio-sorted order out
to $p\le5000$.  The supremum is not attained (see below), so ``$\max$'' should be read as
``$\sup$''.
\end{proposition}

Three cautions, two of which retract earlier readings.

\emph{A search cap of $p<300$ is far too small.}  An earlier version of this table reported
the single rail $\{11\}$ for $7\le q\le10$ and $\{13\}$ for $q=11,12$; every one of those rows
was wrong, the second rail merely sitting between $827$ and $9817$.  The values above come
from branch-and-bound over $p<20\,000$ with a rigorous efficiency bound, cross-checked against
unpruned enumeration.  Likewise \emph{it is false that only $q=2$ supports three rails}: at
$q=3$ the four rails $\{191,193,197,199\}$ are feasible within $p<200$ at cost $0.8276$.  What
is true is only that $q\ge3$ cannot reach three rails \emph{cheaply}.

\emph{Coarsening the division destroys monotonicity of rail cost in $p$, and this is the AC
content of Theorem~\ref{thm:phi} made visible.}  At $q=2$ the cost $\beta_p^{(2)}$ is strictly
decreasing in $p$ --- there are no inversions among primes below $200\,000$, and none can
occur, since $\tfrac{d}{dp}\beta$ has the sign of $\tfrac{\log p}{p+1}-\log\tfrac{2p}{p+1}$,
which is negative for $p\ge3$ because the first term is at most $\tfrac{\log3}{4}=0.275$ and
the second at least $\log\tfrac32=0.406$.  For $q\ge3$ inversions are everywhere: $3271$ of
them at $q=3$ below $200\,000$, $4128$ at $q=4$, $7142$ at $q=12$.  The consequence is
concrete.  At $q=4$,
\[
  \beta_{17}=0.431939,\qquad \beta_{19}=0.453397,\qquad \beta_{23}=0.428556,
\]
so $\{5,17\}$ costs $1.001262$ and $\{5,19\}$ costs $1.022721$ --- both infeasible --- while
$\{5,23\}$ costs $0.997879$ and is feasible.  The optimal second rail is the \emph{largest} of
the three, chosen not by size but by $23\equiv-1\bmod4$ giving the small offset $r=1$.  That
is exactly what a fluctuating registration offset does: it reorders the rails without moving
the ceiling much.

\emph{The maximum at $q=2$ is monotonicity, not rigidity.}  $\lceil p/q\rceil$ is
nonincreasing in $q$, so $\beta_p^{(q)}$ is nondecreasing in $q$ for every fixed $p$
(verified, $3596$ pairs, no violation); the feasible region therefore shrinks monotonically
and $\Lambda$ can only decrease.  Since $q=2$ is the smallest admissible division it is
automatically the maximiser, as it would be for any monotone family.  This is why
Proposition~\ref{prop:ceiling} is not offered as a rigidity.

\emph{No nontrivial optimum is known to be attained.}  Every nondegenerate $\beta_p^{(q)}$ is
irrational, by the argument in \emph{(i)} of the discussion following
Theorem~\ref{thm:niven}.  That alone does \emph{not} settle the
question --- a sum of irrationals can be rational --- and whether
$\sum_{p\in S}\beta_p^{(q)}=1$ is possible for a nondegenerate set $S$ is left open here: the
terms are $\log(p/a)/\log p$ with \emph{different} denominators, so this is not a rational
linear form in logarithms and Baker's theorem does not apply directly.  Numerically the sum
never came within $10^{-11}$ of $1$ over the searched range.  Granting that, each optimum is a
supremum approached through ever longer far tails.  At $q=2$ the tail is exact:
the last admissible rail is $P=398\,328\,266\,023$, contributing
$1/P=2.5104921877\times10^{-12}$ and giving $\Lambda(2)=0.6761904761929866826639\dots$; a
further rail would need $P'\approx1.75\times10^{8602545911616}$.  Coarser divisions have tails
too ($q=3$ needs $P\approx1.5\times10^{54}$, $q=4$ needs $P\approx8\times10^{283}$), so the
far tail is generic and not a feature of the half-turn.  By contrast the spurious optimum
$\{2\}$ at $q\ge3$ had $\sum\beta=1$ \emph{exactly} and \emph{was} attained --- the cleanest
signature that a degenerate rail has been admitted.

The one residue here that is not monotone: $q=2$ is the only division with $\Lambda(q)>1/q$,
by a wide margin ($0.676=1.35\times\tfrac12$, against $\Lambda(q)<1/q$ for every $3\le q\le10$).
This traces to $p=3$ being the half-turn's smallest nondegenerate rail, where the offset
correction is large: $\beta_3^{(2)}=0.369$ against a leading term $\log2/\log3=0.631$.

\section{The window lemma at a $1/q$-turn: the arc fan}

The window lemma behind \#683, \#685 and \#1094 is stated for $p>n/2$ with $k\le n/2$.  Both
halves are the half-turn, and at a $1/q$-turn the single window opens into a fan.

\begin{theorem}[arc fan]\label{thm:fan}
Let $q\ge2$, $n\ge q^{2}$, $1\le k\le n/q$, and let $p>n/q$ be prime.  Put $j=\lfloor n/p\rfloor$,
so $0\le j\le q-1$.  Then
\[
  v_p\binom nk=
  \begin{cases}
    1,&j\ge1\ \text{and}\ jp>n-k,\\
    0,&\text{otherwise.}
  \end{cases}
\]
Equivalently $\{p>n/q:\ p\mid\binom nk\}$ is the disjoint union over $j=1,\dots,q-1$ of the
primes in $\bigl(\tfrac{n-k}{j},\tfrac nj\bigr]$.  Arc $j$ has length exactly $k/j$ and lies
inside the band $(\tfrac n{j+1},\tfrac nj]$.
\end{theorem}

\begin{proof}
$n\ge q^{2}$ gives $p>n/q\ge\sqrt n$, so $p^{2}>n$ and only level $1$ survives; $k\le n/q<p$
gives $\lfloor k/p\rfloor=0$, so $v_p\binom nk=\lfloor n/p\rfloor-\lfloor(n-k)/p\rfloor$, the
number of multiples of $p$ in $(n-k,n]$.  That interval has length $k<p$, so the count is $0$
or $1$, and is $1$ exactly when the largest multiple $jp\le n$ exceeds $n-k$.  Finally
$p>n/q$ forces $j\le q-1$, and $k\le n/q\le n/(j+1)$ places arc $j$ inside its band.
\end{proof}

At $q=2$ only $j=1$ survives and Theorem~\ref{thm:fan} \emph{is} the classical window lemma.
So the single window becomes a fan of $q-1$ arcs, the $j$-th of length $k/j$ at height $n/j$.
The refined window is already turn-free: for $p>\max(k,\sqrt n)$ one has
$p\mid\binom nk\iff(n\bmod p)<k$, with $v_p=1$; it is the fan, not the refinement, that
carries the $q$-dependence.

\begin{theorem}[the ledger is Boolean only at the half-turn]\label{thm:boolean}
For every prime $p>n/q$,
\[
  v_p\binom nk\ \le\ \bigl\lceil\log_2 q\bigr\rceil ,
\]
and the bound is attained.  In particular $v_p\binom nk\in\{0,1\}$ for all $p>n/q$ if and only
if $q\le2$.
\end{theorem}

\begin{proof}
$v_p\ge m$ requires $p^{m}\le n$, while $p>n/q$ gives $n<qp$; hence $p^{m}<qp$, so
$p^{m-1}<q$, so $2^{m-1}<q$ and $m\le\lceil\log_2q\rceil$.  Sharpness: take $p=2$, $k=1$,
$n=2^{\lceil\log_2q\rceil}$.
\end{proof}

\emph{Verified}: the maximum of $v_p\binom nk$ over $p>n/q$, all $n<220$ and $k\le n/2$, is
$1,2,2,3,3,3,3,4,4$ for $q=2,\dots,10$ --- equal to $\lceil\log_2q\rceil$ in every case, with
no violation.

This is the one statement in this note that distinguishes $q=2$ \emph{intrinsically and
qualitatively}: at the half-turn the binomial ledger above the threshold is a single bit, and
at every coarser turn it is a staircase.  Unlike Proposition~\ref{prop:ceiling} it is not an
extremum of a monotone quantity, and unlike Theorem~\ref{thm:niven} it involves no chart.  It
is still a threshold in a monotone family --- $\lceil\log_2q\rceil$ is nondecreasing --- so the
phrase ``unique largest'' is doing the work; what is not automatic is that the value $1$ is
attained at all.  Correspondingly, $v_p\ge2$ with $p>n/q$ forces $n<q^{2}$: at a $1/q$-turn
higher valuations live only in the low range.

The multinomial side is the same lemma one level up.  For $M=A!/(a_1!\cdots a_q!)$ with
$A=\sum_i a_i$, the level-$j$ entry $\bigl\lfloor\bigl(\sum_i(a_i\bmod p^{j})\bigr)/p^{j}\bigr\rfloor$
is the \emph{wrap number} of the $q$ residues on the circle $\mathbb Z/p^{j}$, valued in
$\{0,\dots,q-1\}$; and for $p>\max(\max_i a_i,\sqrt A)$ one gets $v_p(M)=\lfloor A/p\rfloor$,
taking the value $m$ exactly on $p\in\bigl(\tfrac A{m+1},\tfrac Am\bigr]$.  At $q=2$ with
$a=(k,n-k)$ this is $v_p=\lfloor n/p\rfloor\in\{0,1\}$, the classical window lemma again: the
window \emph{value} $1$ becomes a staircase running through every value $0,\dots,q-1$.

\section{Erd\H{o}s \#376 at general $q$}

\#376 asks whether $\binom{2n}{n}$ is coprime to $105=3\cdot5\cdot7$ infinitely often;
equivalently, by Corollary~\ref{cor:carrier}, whether infinitely many $n$ lie in
$G_3^{(2)}\cap G_5^{(2)}\cap G_7^{(2)}$.  The $q$-analogue collapses immediately.

\begin{proposition}\label{prop:376}
For every $q\ge3$ the only $n\ge0$ with $M_q(n)$ coprime to $105$ is $n=0$.
\end{proposition}

\begin{proof}
$3\le q$, so $\lceil3/q\rceil=1$ and the only admissible base-$3$ digit is $0$; hence
$G_3^{(q)}=\{0\}$ and $3\mid M_q(n)$ for every $n\ge1$.
\end{proof}

So the three-rail problem is open at $q=2$ and refuted for every $q\ge3$, with no intermediate
regime: there is no coarser division at which \#376 is easier, because the smallest rail dies
first.  The per-rail dimensions $d_p^{(q)}=1-\beta_p^{(q)}=\log\lceil p/q\rceil/\log p$ make
this quantitative.  At $q=2$ they are $0.630930,\ 0.682606,\ 0.712414$ summing to $2.025950$;
at $q=3$, $0,\ 0.430677,\ 0.564575$ summing to $0.995252$; at $q=4$, $0.786884$; at $q=5,6$,
$0.356207$; at $q\ge7$, zero.

Two thresholds are easy to conflate here and should not be.  The condition $\sum_p d_p>1$ is
the \emph{two}-rail criterion; three simultaneous digit constraints add codimensions, so the
right test is $\sum_p d_p>2$, equivalently $\sum_p\beta_p<1$.  Only $q=2$ passes it, with
predicted dimension $\sum d_p-2=0.025950$.  The $q=3$ value $0.995252$ is \emph{not} a near
miss of the three-rail test --- it misses by $1.005$.  Note also the bookkeeping identity
$1-\sum_{p\in\{3,5,7\}}\beta_p^{(2)}=\sum_p d_p-2$: the leftover budget in
Proposition~\ref{prop:ceiling} and the predicted \#376 dimension are the same number,
$0.025950322273$.

\emph{Verified}: $|G_p^{(q)}\cap[0,p^{J})|=\lceil p/q\rceil^{J}$ exactly, for
$p\in\{3,5,7,11\}$, $2\le q\le8$, $1\le J\le4$ --- no failure.  So the per-rail box count is
exact at every scale and nothing heuristic enters below the multi-rail level; the budget rule
``$\sum\beta<1\Rightarrow$ infinitely many'' is precisely the assumption that the digit
constraints at distinct bases are independent, and that assumption is the whole content of
\#376.

\begin{remark}[attribution of the budget rule]\label{rem:attrib}
The rule ``$\sum_p\beta_p^{(q)}<1\Rightarrow$ infinitely many'' is not ours.  With
$\kappa_j=1/q$ it is verbatim Conjecture~1 of Bloom and Croot, \emph{Integers with small digits
in multiple bases}, arXiv:2509.02835: for bases $g_j$ with associated $\kappa_j\in(0,1]$,
if $\sum_j\log_{g_j}\bigl(g_j/\lceil\kappa_jg_j\rceil\bigr)<1$ then there are infinitely many
$n$ all of whose base-$g_j$ digits are $<\kappa_jg_j$ --- and $\lceil\kappa_jg_j\rceil$ is
exactly our $\lceil p/q\rceil$.  They compute the Graham case
$\log_3\tfrac32+\log_5\tfrac53+\log_7\tfrac74=0.974\dots$, our $0.974050$.  The finiteness
clause and the predicted dimension $\sum_pd_p-(r-1)$ are Conjecture~1.4 of Han Yu,
\emph{Fractal projections with an application in number theory}, arXiv:2004.05924, stated there
for $\kappa=\tfrac12$.  The exponent $0.02595$ is Pomerance's published heuristic (Amer.\ Math.\
Monthly \textbf{122} (2015), \S4), recorded on OEIS A030979.  What the $q$-family contributes is
only the dictionary $\kappa=1/q\leftrightarrow(qn)!/(n!)^{q}$, connecting the digit side to a
multinomial; that is a dictionary entry, not a theorem.  Note also where
Remark~\ref{rem:q3crit} sits: Bloom--Croot have general $\kappa$ but only the infinitude
clause, Han Yu has the finiteness clause but only at $\kappa=\tfrac12$, and the $q=3$ two-rail
case needs both at once --- it is the uncovered corner between two printed conjectures.
\end{remark}

\begin{remark}[a sharper critical case than \#376 itself]\label{rem:q3crit}
At $q=3$ the two smallest usable rails give
$\beta_5^{(3)}+\beta_7^{(3)}=1.004748407873$, missing the budget by $0.475\%$ --- the thinnest
margin anywhere in the family.  There is a sharper way to say this.  The only unconditional
two-base theorem in the literature is Theorem~1 of Erd\H{o}s, Graham, Ruzsa and Straus,
\emph{On the prime factors of $\binom{2n}{n}$}, Math.\ Comp.\ \textbf{29} (1975), 83--92, p.~84:
if $\tfrac{A}{p-1}+\tfrac{B}{q-1}>1$ then infinitely many integers have all base-$p$ digits
$<A$ and all base-$q$ digits $<B$.  Our case is base $5$ with digits $<5/3$, i.e.\ $A=2$, and
base $7$ with digits $<7/3$, i.e.\ $B=3$, giving
\[
  \frac{A}{p-1}+\frac{B}{q-1}=\frac24+\frac36=\frac12+\frac12=1
\]
\emph{exactly} --- failing the strict inequality by zero.  So this is not merely a near miss of
a heuristic budget; it is the exact boundary case of the one theorem that would settle it.
(Bloom and Croot restate EGRS Theorem~1 in the inequivalent form
$\tfrac{\lceil\kappa_1g_1\rceil-1}{g_1-1}+\tfrac{\lceil\kappa_2g_2\rceil-1}{g_2-1}\ge1$, under
which our case gives $\tfrac14+\tfrac26=0.583$; the case is uncovered either way, but the two
forms must not be quoted interchangeably.)  The corresponding question, whether $(3n)!/(n!)^{3}$ is coprime
to $35$ infinitely often, is therefore predicted \emph{finite}, and barely.  A search to
$2.98\times10^{17}$ finds exactly nine solutions,
\[
  0,\ 1,\ 750,\ 751,\ 3150,\ 3151,\ 4761974610150,\ 4761974610155,\ 4761974610156,
\]
with a desert of nine orders of magnitude between $3151$ and $4.76\times10^{12}$ that then
ends.  Near-critical verdicts of this kind are not numerically decidable --- at $q=4$ the pair
$\{5,23\}$ has predicted dimension $0.00212$ and its solution count went from $9$ to $30$
between $10^{9}$ and $3\times10^{17}$, in tight clusters, where the prediction allows a factor
$1.04$ --- so this is a prediction on record, not a result.
\end{remark}

\begin{remark}[do not rewrite the digit criterion with a floor]\label{rem:floor}
Corollary~\ref{cor:carrier} must be kept in its strict form, ``every base-$p$ digit
$<p/q$''.  The tempting rewrite ``every digit $\le\lfloor p/q\rfloor$'' is false exactly when
$p=q$: there $\lfloor p/q\rfloor=1$ but $p/q=1$, so the strict form correctly forces every
digit to $0$, matching $G_p^{(p)}=\{0\}$, while the floor form wrongly admits the digit $1$.
At $p=q=2$ the floor version would assert that $\binom{2n}{n}$ is always odd.  The strict form
is valid for \emph{all} primes with no side condition, and the alphabet size $\lceil p/q\rceil$
is likewise universally correct, so $\beta_p^{(q)}$ has no exceptional regime.
\end{remark}

\begin{remark}[how much of this is really rigidity]\label{rem:deflate}
The three statements above all return $q\in\{1,2\}$, and it would be easy to read that
agreement as threefold confirmation.  It is not.  Proposition~\ref{prop:ceiling} is
monotonicity in $q$ and would hold for any monotone family; Theorem~\ref{thm:niven} is true
with nil content and its cell is chart-created.  Only Theorem~\ref{thm:phi}(AC) is neither an
extremum nor a convention.

Are (AC) and Niven the same statement twice?  \emph{No, and the near miss is instructive.}
Niven's theorem \emph{is} a $\varphi$-condition --- $2\cos(2\pi k/n)$ generates
$\mathbb Q(\zeta_n)^{+}$, of degree $\varphi(n)/2$, so rationality is $\varphi(n)\le2$, i.e.\
$n\in\{1,2,3,4,6\}$ --- but on a \emph{different modulus}.  Niven's modulus is the order
$n=6$ of the root of unity attached to $\cos\theta=\tfrac12$, with $\varphi(6)=2$: an
\emph{index} condition, $(\mathbb Z/n)^{\times}=\{\pm1\}$.  Ours is $q=2$ with
$\varphi(2)=1$: an \emph{order} condition, $(\mathbb Z/q)^{\times}=\{1\}$.  That the two cut
out the same set $\{1,2\}$ is extensional.  Perturbing the family separates them: with
$\delta=1/(q-1)$ Niven selects $\{2,3\}$ and $\varphi$ selects $\{2\}$; with
$\delta=(q-1)/(q+1)$ Niven selects $\{3\}$ and $\varphi$ selects $\{2\}$ --- disjoint.

They do share an ancestor, the involution $x\mapsto-x$: $\varphi(q)=1$ says $-1\equiv1\bmod q$,
and $\varphi(n)\le2$ says $(\mathbb Z/n)^{\times}=\{\pm1\}$.  But ours acts on the digits, and
is therefore intrinsic to the object; Niven's acts on the cyclotomic field of the chart's
angle, and is not.  Only one of the two ancestors is present in the construction.

The claim worth carrying forward is therefore the narrow one: the mean registration offset is
$\tfrac1{2p}$ at every division, and the half-turn is the unique division at which that offset
does not fluctuate.
\end{remark}

\begin{remark}[harmonic content]
The object's own scale is the circle cut into $q$ sectors, never unit $1$.  The DC mode is the
sector index $\lfloor q\{n/p^{j}\}\rfloor$, evaluated exactly in Theorem~\ref{thm:ledger} with
no remainder; the entire AC content at a prime above the threshold is the single sector
occupied by $\rho/p$.  The difficulty lives on the carrier, in the smooth window that
Theorem~\ref{thm:anti}(a) forces, and Theorem~\ref{thm:qge} shows that difficulty is
\emph{purchased by the coarseness of the division}: at $q\ge k+1$ the carry budget of $M_q$
covers the window outright.
\end{remark}

\begin{remark}[register]
\textbf{Proved}: Theorem~\ref{thm:ledger}, Corollary~\ref{cor:carrier},
Theorem~\ref{thm:anti}, Corollary~\ref{cor:crit}, Theorem~\ref{thm:qge},
Theorem~\ref{thm:phi}, Theorem~\ref{thm:niven}.  All are elementary from Legendre, Kummer,
Dirichlet and Niven, and every one was checked numerically before being written.
\textbf{Measured}: Proposition~\ref{prop:ceiling} and the least-$n$ table, both within the
stated search ranges.
\textbf{Not proved}: Erd\H{o}s \#396 itself, i.e.\ the case $q=2$ for every $k$; and, for each
fixed $q\ge3$, the analogue for every $k$ --- Theorem~\ref{thm:qge} covers only $k<q$.
\textbf{Falsifier}: a prime $p>\max(qk,\sqrt{qn})$ dividing some $n-i$ with
$v_p(M_q(n))>0$ would refute Theorem~\ref{thm:anti}(a); none was found in $66\,896$ tested
incidences across five values of $q$.  A division $q\ge3$ with constant registration offset
would refute Theorem~\ref{thm:phi}; none exists, by $\varphi(q)\ge2$.
\end{remark}
\end{document}
